% ----------------------
\subsection{Section 131 (16--17 May 1843)}
% ----------------------
	\label{DC131}
	
	% -----
	\subsubsection{Comment on this section}
	% -----
	
		Be careful with reading this section of the DC. The original sources, discussed below, have much more information and in some cases are meaningfully different from the account we have in this section. Best to read the originals and leave this section alone.
	
	% -----
	\subsubsection{The original source for DC 131}
	% -----
	
		This section is a compilation of things JS said on a few different occasions over two different days, on 16 and 17 May 1843. The three different accounts are some \hyperref[EMS-1843-JS-16-MAY-1843-WC]{instructions} recorded by William Clayton on 16 May, a \hyperref[EMS-1843-JS-17-MAY-1843-A-WC]{discourse} recorded by William Clayton on 17 May, and another \hyperref[EMS-1843-JS-17-MAY-1843-B-WC]{discourse} recorded by William Clayton later on 17 May.
		
		See the following list for a breakdown of where this material comes from.
		
		\begin{compactdesc}
			\item[Verse 1] \hyperlink{EMS-1843-JS-16-MAY-1843-WC-d}{Instruction, 16 May 1843, passage d}
			\item[Verse 2] \hyperlink{EMS-1843-JS-16-MAY-1843-WC-d}{Instruction, 16 May 1843, passage d}
			\item[Verse 3] \hyperlink{EMS-1843-JS-16-MAY-1843-WC-d}{Instruction, 16 May 1843, passage d}
			\item[Verse 4] \hyperlink{EMS-1843-JS-16-MAY-1843-WC-d}{Instruction, 16 May 1843, passage d}
			\item[Verse 5] \hyperlink{EMS-1843-JS-17-MAY-1843-A-WC-b}{Discourse, 17 May 1843--A, passage b}
			\item[Verse 6] \hyperlink{EMS-1843-JS-17-MAY-1843-A-WC-b}{Discourse, 17 May 1843--A, passage b}
			\item[Verse 7] \hyperlink{EMS-1843-JS-17-MAY-1843-B-WC-b}{Discourse, 17 May 1843--B, passage b}
			\item[Verse 8] \hyperlink{EMS-1843-JS-17-MAY-1843-B-WC-b}{Discourse, 17 May 1843--B, passage b}
		\end{compactdesc}
	
	% -----
	\subsubsection{Historical Background}
	% -----
		\label{DC131-HB}
		
		% --
		\paragraph{JSP}
		% --
		
			%
			\subparagraph{Instruction, 16 May 1843}
			%
		
				``On 16 May 1843 at the residence of Benjamin F. and Melissa LeBaron Johnson in Macedonia, Illinois, JS instructed the Johnsons and William Clayton `on the priesthood.' Benjamin Johnson recalled that `in the evening,' JS invited him and Melissa `to come and sit down, for he wished to marry us according to the law of the Lord. I thought it a joke, and said I should not marry my wife again, unless she courted \uline{me}, for I did it all, the first time. He chided my levity, told me he was in earnest, and so it proved; for we stood up and were sealed by the Holy Spirit of promise.' Contemporaneous records, however, show that when he later recounted the story, Johnson conflated when he first heard this teaching and when he and Melissa were actually sealed. According to William Clayton's journal, the couple was `united in an everlasting covenant' on 20 October 1843.

				``JS taught that if a man and a woman were sealed as husband and wife by the power of the priesthood, their marriage would continue beyond death, and they would live together in marriage in the highest level of heaven. His teachings included the idea that couples properly sealed would have `children in the resurrection.' Based on JS's statement to Johnson that he knew in whom he could confide, it seems he expected the Johnsons not to discuss this doctrine with others. 

				``William Clayton was present during JS's instruction and took notes. However, in view of the uniformity of the ink flow and character formation in his journal, it appears that Clayton created the entries covering 14 May through most of 19 May in one sitting. It is likely the entries were based upon earlier notes that are no longer extant.''
				
			%
			\subparagraph{Discourse, 17 May 1843--A}
			%
			
				``On 17 May 1843, JS preached a discourse in Macedonia, Illinois, on the first chapter of 2 Peter in the New Testament. In the discourse, JS revisited a subject he had addressed three days before in Morley Settlement, Illinois---the individual Christian believer's obtaining assurance of having eternal life. As he did in the earlier discourse, JS emphasized the importance of knowledge in this process and explained the `more sure word of prophecy' mentioned in 2 Peter 1:19, which he interpreted as receiving divine confirmation of being `sealed up unto eternal life.' In this 17 May discourse, JS indicated that this divine affirmation would come through the priesthood. 

				``Having been apprised on the evening of 16 May that JS planned to preach the following day, the Saints congregated at 10:00 a.m. on 17 May, possibly in a schoolhouse built by church members in Macedonia. Samuel Prior, a Methodist preacher who was visiting Macedonia to observe the Latter-day Saints in general and JS in particular, attended the meeting and subsequently wrote an account of his experience, which was published in the \textit{Times and Seasons}. Prior described taking his `seat in a conspicuous place in the congregation' and waiting for JS's arrival, after which JS `commenced calmly and continued dispassionately to pursue his subject. . . . He glided along through a very interesting and elaborate discourse, with all the care and happy facility of one who was well aware of his important station, and his duty to God and man, and evidencing to me, that he was well worthy to be styled ``\textit{a workman rightly dividing the word of truth},'' and giving without reserve, ``\textit{saint and sinner his portion in due season}.''' Prior stated, `I was compelled to go away with a very different opinion from what I had entertained when I first took my seat to hear him [JS] preach.'

				``William Clayton, who accompanied JS to Macedonia, wrote an account of the dis-course in his journal. The relatively polished nature of the featured text suggests that Clayton reconstructed JS's words after the fact, likely from notes taken at the time. Based on the uniformity of the ink flow and character formation in his journal, it appears that Clayton created the entries covering 14 May through most of 19 May in one sitting. The original notes are apparently not extant.''
				
			%
			\subparagraph{Discourse, 17 May 1843--B}
			%
			
				``On the evening of 17 May 1843, JS preached a discourse in Macedonia, Illinois, in response to a sermon given by Samuel Prior, a visiting Methodist preacher, who had been `invited to preach' earlier that evening before a `large and respectable' congregation of Latter-day Saints. Prior wrote an article that was published in the \textit{Times and Seasons} shortly after JS gave the sermon featured here. After Prior delivered his unrecorded discourse, JS sought to correct Prior's interpretation of Genesis 2 and the `eternal duration of matter'---doing so `mildly, politely, and affectingly,' according to Prior.

				``During his remarks, JS questioned the King James Bible's language in Genesis 2:7, which states, `The Lord God formed man of the dust of the ground, and breathed into his nostrils the breath of life; and man became a living soul.' According to Clayton's notes of the discourse, JS also discussed aspects of Hebrew in his argument. JS argued that the `breath of life' referred to Adam's spirit, and he included comments on the Hebrew word \textit{ruach} (breath or spirit), arguing that when the word is applied to Eve, \textit{life} should be rendered \textit{lives}.

				``Genesis 2:7 prompted pronouncements from religious thinkers of JS's era on the nature of bodies and souls and on their duration throughout eternity. Methodist minister Adam Clarke, for instance, wrote that `man is a \textit{compound} being, having a body and a soul, distinctly and separately created; the body out of the dust of the earth, the soul immediately breathed from God himself.' A widely circulated nineteenth-century theological dictionary stated that the `generally received opinion' held that the human soul `began to exist in his mother's womb.' Though it is impossible to reconstruct Prior's sermon, JS's response makes clear that Prior used his sermon to engage this topic. JS responded to Prior and offered his own thoughts on the nature of the human spirit. 

				``JS's own teachings about the soul and matter differed from those of other Christians of his day in some important respects. Common teachings stated that the soul consisted of `the vital, immaterial, active substance, or principle, in man.' For some years, JS had taught that all beings were composed of eternal matter. A May 1833 revelation stated that `the Elements are eternal.'%
					\footnote{%
						% \label{}%
						See \hyperref[DC93-33]{DC 93:33}.
						}
				On 30 August 1840, JS preached a `discourse on Eternal Judgement and the Eternal Duration of matter.'%
					\footnote{%
						% \label{}%
						See \hyperref[EMS-1840-JS-30-AUG-1840]{this discourse}; there isn't much more there.
						}
				In 1842, JS published an editorial directly opposing the traditional idea that the spirit was immaterial. He argued, much as he did in the discourse featured here, that `the body is supposed to be organized matter, and the spirit by many is thought to be immaterial, without substance. With this latter statement we should beg leave to differ---and state that spirit is a substance; that it is material, but that it is more pure, elastic, and refined matter than the body.' Also in 1842, the published Book of Abraham stated that `the Gods formed man from the dust of the ground, \textit{and took his spirit, that is the man's spirit, and put it into him}, and breathed into his nostrils the breath of life, and man became a living soul.' The discourse featured here shares the Book of Abraham's emphasis on the eternal nature of humankind and the idea that spirits were embodied at birth. 

				``William Clayton was present during JS's discourse and took notes. However, based on the uniformity of the ink flow and character formation in his journal, it appears that Clayton created the entries covering 14 May through most of 19 May in one sitting. It is likely the entries were based upon earlier, nonextant notes.''
		
	% -----
	\subsubsection{Versions}
	% -----
		\label{DC131-V}
		
		See the \href{https://drive.google.com/file/d/1goAvNz8jS4R8slYfMG_db55Ty2z_vmgK/view?usp=sharing}{sections comparison} document.
	
		\begin{compactdesc}
			\item[JSP] \hyperref[EMS-1843-JS-2-APR-1843]{2 April 1843}
			\item[BC] ---
			\item[DC 1835] ---
			\item[DC 1844] ---
			\item[DN] 6:28, 29 (17 and 24 September 1856), 217, 225
			\item[MS] 21:7, 9 (12 and 26 February 1859), 108, 142
		\end{compactdesc}

	% -----
	\subsubsection{Section 131:1} % *DC131-1 
	% -----
		\label{DC131-1}

			\footnote{%
				% \label{}%
				Prefacing this record of the remarks that we have as verses 1--4 of this section, William Clayton wrote, ``Before we retired the prest. gave bro Johnson \& wife some instructions on the priesthood.''
				}%
		IN the celestial glory there are three heavens or
			\footnote{%
				% \label{}%
				``Degree'' is pretty clearly a Masonic word. And you progress by degrees or through the degrees, so I think JS probably intended it to be understood as a progression from one degree to the other, occurring either here on Earth in some way, or after the resurrection.
				}%
		\hyperref[DEGREE]{degrees};

		% \begin{framed}
		% \scriptsize
			% \begin{compactdesc}
				% \item[JSP] 
				% % \item[BC] 
				% % \item[]
			% \end{compactdesc}
		% \end{framed}

	% -----
	\subsubsection{Section 131:2} % *DC131-2 
	% -----
		\label{DC131-2}

		And in order to obtain the highest, a man must enter into this order of the priesthood [meaning%
			\footnote{%
				% \label{}%
				In the original source, they have been talking about marriage for eternity, and so the words ``this order of the priesthood'' make a clearer reference than they do here, which is why you get this bracketed explanation. See the earlier \hyperlink{EMS-1843-JS-16-MAY-1843-WC-b}{passage b} of the original source.
				}
		the new and everlasting covenant of marriage];

		% \begin{framed}
		% \scriptsize
			% \begin{compactdesc}
				% \item[JSP] 
				% % \item[BC] 
				% % \item[]
			% \end{compactdesc}
		% \end{framed}

	% -----
	\subsubsection{Section 131:3} % *DC131-3 
	% -----
		\label{DC131-3}

		And if he does not, he cannot obtain it.

		% \begin{framed}
		% \scriptsize
			% \begin{compactdesc}
				% \item[JSP] 
				% % \item[BC] 
				% % \item[]
			% \end{compactdesc}
		% \end{framed}

	% -----
	\subsubsection{Section 131:4} % *DC131-4 
	% -----
		\label{DC131-4}

		He may enter into the other,%
			\footnote{%
				% \label{}%
				The other order of the priesthood, presumably; earlier he was talking about \textit{obtaining} the highest degree, not \textit{entering into} it, so I think it's talking here about entering into some other order of the priesthood.
				}
		but that is the end of his kingdom; he cannot have an increase.%
			\footnote{%
				% \label{}%
				The original source, at \hyperlink{EMS-1843-JS-16-MAY-1843-WC-b}{passage b}, also has more here---note that ``increase'' is spoken of separately from ``children'': ``those who are married by the power \& authority of the priesthood in this life \& continue without committing the sin against the Holy Ghost will continue to increase \& have children in the celestial glory.''
				}

		% \begin{framed}
		% \scriptsize
			% \begin{compactdesc}
				% \item[JSP] 
				% % \item[BC] 
				% % \item[]
			% \end{compactdesc}
		% \end{framed}

	% -----
	\subsubsection{Section 131:5} % *DC131-5 
	% -----
		\label{DC131-5}

		(May 17th, 1843.)
			\footnote{%
				% \label{}%
				This is commentary on \hyperref[2PETER-1-19]{2 Peter 1:19}. Check the original source \hyperref[EMS-1843-JS-17-MAY-1843-A-WC]{here}, which has more useful stuff.
				}%
		The more sure word of prophecy means a man's knowing that he is sealed up unto eternal life, by revelation and the spirit of prophecy, through the power of the Holy Priesthood.

		% \begin{framed}
		% \scriptsize
			% \begin{compactdesc}
				% \item[JSP] 
				% % \item[BC] 
				% % \item[]
			% \end{compactdesc}
		% \end{framed}

	% -----
	\subsubsection{Section 131:6} % *DC131-6 
	% -----
		\label{DC131-6}

		It is impossible for a man to be saved in ignorance.

		% \begin{framed}
		% \scriptsize
			% \begin{compactdesc}
				% \item[JSP] 
				% % \item[BC] 
				% % \item[]
			% \end{compactdesc}
		% \end{framed}

	% -----
	\subsubsection{Section 131:7} % *DC131-7 
	% -----
		\label{DC131-7}

		There is no such thing as immaterial matter.  All spirit is matter, but it is more fine or pure, and can only be discerned by purer eyes;

		% \begin{framed}
		% \scriptsize
			% \begin{compactdesc}
				% \item[JSP] 
				% % \item[BC] 
				% % \item[]
			% \end{compactdesc}
		% \end{framed}

	% -----
	\subsubsection{Section 131:8} % *DC131-8 
	% -----
		\label{DC131-8}

		We cannot see it; but when our bodies are purified we shall see that it is all matter.

		% \begin{framed}
		% \scriptsize
			% \begin{compactdesc}
				% \item[JSP] 
				% % \item[BC] 
				% % \item[]
			% \end{compactdesc}
		% \end{framed}